% ============================================================
%  CONCLUSION GÉNÉRALE
% ============================================================
\chapter*{Conclusion Générale}
\addcontentsline{toc}{chapter}{Conclusion Générale}
\markboth{Conclusion Générale}{Conclusion Générale}

\section*{Bilan du projet}

Ce projet de fin d'études a abouti à la réalisation d'un système complet de maintenance
prédictive avec intelligence artificielle, baptisé \textbf{PC Technician Assistant}.
Le système répond à la problématique initiale en transformant la maintenance réactive
en maintenance proactive grâce à l'analyse automatisée des métriques matérielles.

Les principaux résultats obtenus sont les suivants :

\begin{itemize}
  \item \textbf{Architecture microservices fonctionnelle} : quatre services indépendants
        (agent Python, backend Node.js, service ML Flask, frontend React) communiquant
        via des API REST sécurisées par JWT et tokens Bearer ;
  \item \textbf{Collecte automatisée} : 7,8 millions d'enregistrements de métriques
        accumulés sur 20 machines, avec collecte horaire sans interruption via
        \texttt{psutil} et \texttt{smartctl} ;
  \item \textbf{Dashboard avec double prédiction ML} : le tableau de bord affiche
        côte à côte la \textbf{probabilité de panne Random Forest à 30 jours} (colonne
        ``Proba. Panne (RF)'') et la \textbf{santé disque LSTM} (colonne ``Santé Disque
        (LSTM)'') pour chaque machine, offrant deux perspectives complémentaires sur le
        risque de défaillance ;
  \item \textbf{Pipeline ML opérationnel} : extraction de 63 caractéristiques statistiques
        et génération de prédictions à 7/14/30 jours avec versioning automatique des
        modèles (7 versions sauvegardées, activation conditionnelle à +5\% d'accuracy) ;
  \item \textbf{LSTM entraîné sur données réelles} : modèle PyTorch entraîné sur le
        dataset Backblaze Q1 2020 (874 892 lignes, 4 993 disques), atteignant un rappel
        de 100\% sur les pannes réelles et un ROC AUC de 1,000 ;
  \item \textbf{Chatbot hybride} : assistant conversationnel combinant base de
        connaissances déterministe, requêtes SQL live et LLM local (Ollama/LLaMA 3.2:1b),
        évalué avec ROUGE-1 = 0,4548 et ROUGE-2 = 0,3043 sur 20 questions de référence ;
  \item \textbf{Sécurité multicouche} : JWT (24h), bcrypt (10 rounds), RBAC
        Admin/Technicien, protection injection SQL via Sequelize ORM.
\end{itemize}

\section*{Difficultés rencontrées}

\textbf{Labels synthétiques pour le Random Forest :} Le principal défi technique a été
l'absence de données réelles de pannes pour l'entraînement du modèle Random Forest,
nécessitant le recours à des labels synthétiques basés sur des seuils métier (CPU $> 80\%$,
RAM $> 85\%$, disque $> 90\%$). Cette limitation impacte la précision du modèle (50--70\%),
qui s'améliorerait significativement avec des données de pannes historiques réelles.

\textbf{Disponibilité de SMART sur les machines de test :} La lecture des données SMART
via \texttt{smartctl} n'est pas disponible sur toutes les machines (machines virtuelles,
conteneurs Docker). Le module \texttt{SmartReader} a dû implémenter une stratégie de
fallback robuste pour garantir le fonctionnement de l'agent dans tous les environnements.

\textbf{Déséquilibre de classes dans le LSTM :} Le dataset Backblaze est fortement
déséquilibré (moins de 1\% de pannes). La combinaison d'oversampling, de
\texttt{pos\_weight = 50} et de \texttt{StepLR} a permis d'atteindre un rappel de 100\%,
mais au prix d'une précision de 92\% (8\% de fausses alarmes acceptables).

\section*{Perspectives d'évolution}

Plusieurs axes d'amélioration ont été identifiés pour les versions futures du système :

\begin{enumerate}
  \item \textbf{Données réelles de pannes} : collecter et intégrer des logs de
        défaillances réels pour remplacer les labels synthétiques et améliorer la
        précision du Random Forest ;
  \item \textbf{Intégration du LSTM au dashboard} : le modèle LSTM est désormais
        intégré au tableau de bord avec une colonne dédiée ``Santé Disque (LSTM)''
        dans le composant \texttt{MachineList}, affichée aux côtés de la colonne
        ``Proba. Panne (RF)'' du Random Forest ;
  \item \textbf{Chatbot RAG} : implémenter une architecture RAG
        (Retrieval-Augmented Generation) pour améliorer la qualité des réponses sur
        des questions spécifiques à l'infrastructure, en indexant la documentation
        technique dans une base vectorielle ;
  \item \textbf{Module d'interventions} : développer un module de gestion des
        interventions de maintenance avec historique, planification et suivi des
        techniciens ;
  \item \textbf{Export PDF} : générer des rapports de maintenance automatiques au
        format PDF, incluant les graphiques de tendances et les prédictions ;
  \item \textbf{Application mobile} : développer une application mobile React Native
        pour les techniciens terrain, avec notifications push pour les alertes CRITICAL ;
  \item \textbf{Déploiement Kubernetes} : migrer vers une infrastructure Kubernetes
        pour la scalabilité à grande échelle et la haute disponibilité.
\end{enumerate}

\section*{Apport académique et professionnel}

Ce projet a permis de mettre en pratique et d'approfondir des compétences dans plusieurs
domaines : développement full-stack (React, Node.js, Python), apprentissage automatique
en production (Random Forest, LSTM, Isolation Forest), traitement des séries temporelles
(features statistiques glissantes, tendances linéaires), intégration de LLM locaux
(Ollama), évaluation NLP (métriques ROUGE), et déploiement conteneurisé (Docker).

Il constitue une démonstration concrète de l'application de l'intelligence artificielle
à un problème industriel réel, en intégrant l'ensemble du cycle de vie d'un système ML :
de la collecte des données brutes jusqu'à l'interface utilisateur finale, en passant par
l'ingénierie des features, l'entraînement, l'évaluation et le déploiement en production.
